% Standalone, Overleaf-ready course report template.
% GECCO19 uses ACM SigConf. nonacm suppresses irrelevant publication metadata.
% One source file, inline bibliography, no external figure files required yet.
% Person 3 owns this file; Persons 1/2 verify their Methods descriptions.
\documentclass[sigconf,nonacm]{acmart}
\settopmatter{printacmref=false,printfolios=true}
\renewcommand\footnotetextcopyrightpermission[1]{}
\usepackage{amsmath}
\usepackage{booktabs}
\usepackage{graphicx}
\usepackage{xcolor}
\newcommand{\todo}[1]{\par\smallskip\noindent\textcolor{blue!60!black}{\textbf{TO COMPLETE:} #1}\par\smallskip}
\newcommand{\figureplaceholder}[2]{%
  \IfFileExists{#1}{\includegraphics[width=\linewidth]{#1}}{%
    \fbox{\parbox[c][30mm][c]{0.92\linewidth}{\centering #2\\[3mm]
    Upload the final figure here. No experimental data have been inserted.}}}}
\title{Evolving Controller Topology for Targeted Robot Locomotion}
\author{Team name / number: TO COMPLETE}
\affiliation{\institution{Evolutionary Computing, 2026}\country{The Netherlands}}
\renewcommand{\shortauthors}{Team TO COMPLETE}

\begin{document}
\onecolumn
\thispagestyle{empty}
\begin{center}
\vspace*{20mm}
{\LARGE Evolutionary Computing\par}
\vspace{6mm}
{\Large Standard Assignment -- Task II\par}
\vspace{12mm}
{\huge\bfseries Evolving Controller Topology\\for Targeted Robot Locomotion\par}
\vspace{15mm}
{\large Team number / name: TO COMPLETE\par}
\vspace{8mm}
\begin{tabular}{ll}
\toprule
Team member & Student ID\\
\midrule
Person 1: TO COMPLETE & TO COMPLETE\\
Person 2: TO COMPLETE & TO COMPLETE\\
Person 3: TO COMPLETE & TO COMPLETE\\
% Add a row for every other registered team member, if applicable.
\bottomrule
\end{tabular}
\vspace{12mm}
\par Date: TO COMPLETE\par
\vfill
\par Draft report template; experiments have not yet been run.\par
\end{center}
\clearpage
\twocolumn
\setcounter{page}{1}
\maketitle

% Writing budget: Introduction ~0.75 page; Methods ~2 pages;
% Results/Discussion ~2.5 pages; Conclusions ~0.25 page.
% Hard limit: 6 CONTENT pages. Cover and bibliography do not count.
% All proposed methods below MUST be updated to match the final implementation.

\section{Introduction}
Neuroevolution optimizes neural-network controllers using their performance on a
task, rather than a collection of desired output examples. A conventional
approach fixes the network topology and evolves its connection weights.
Topology-and-weight evolution additionally allows the network structure to
change during search. NEAT motivates starting with small networks and adding
complexity incrementally, alongside other mechanisms including speciation and
historical markings~\cite{stanley2002}.

This study investigates a restricted version of incremental complexification:
adding hidden neurons to a single-layer controller while evolving its weights.
It isolates growth without implementing full NEAT. The task is targeted
locomotion of a fixed robot body in a fixed simulated environment.

\paragraph{Research question.}
Does incrementally adding hidden neurons improve target-reaching performance
compared with fixed-topology weight evolution, at the same simulation-evaluation
budget?

\paragraph{Hypothesis.}
Starting with a small search space and introducing additional capacity gradually
may improve performance within a limited budget. Alternatively, growth may
consume search effort without improving locomotion, or a larger fixed network
may obtain comparable performance without structural mutation.

\todo{Person 3: explain why the comparison matters using the course lecture.
Briefly distinguish the small fixed control from the larger fixed control.
Document the instructor's scope clarification when obtained. Treat the
hypothesis as a prediction, not a result.}

\section{Methods}
\subsection{Task, robot and simulation}
\todo{Person 2 supplies facts; Person 3 writes: identify the exact John Set body
and environment, ARIEL version/commit, Python and MuJoCo versions, timestep,
spawn, target, duration, action application policy, observation normalization,
and how simulation state and the global callback are reset between individuals.}

The proposed setup uses \texttt{john\_set.gecko()} in
\texttt{SimpleFlatWorld}, spawned at $(0,0,0.1)$ with target ground-plane
coordinates $(2,0)$. Each final evaluation lasts 15 simulated seconds.
These choices are provisional until the implementation and pilot are verified.
All conditions use the same task, body, environment and simulation settings.

The fitness of a controller is the remaining horizontal distance
\begin{equation}
 d = \sqrt{(x_{\mathrm{final}}-x_{\mathrm{target}})^2
           +(y_{\mathrm{final}}-y_{\mathrm{target}})^2},
\end{equation}
which is minimized. Height is excluded, and no network-complexity penalty is
included in the fitness function.
\todo{Specify any numerical-failure score and how failures count against the
budget. Do not replace failed runs silently or hide coding errors as bad fitness.}

\subsection{Controller and representation}
\todo{Person 2: list every observation, its order/scale, and the final input and
output dimensions. Include target information and periodic time features.}

The planned controller is a dense, single-hidden-layer network with
\(\mathbf{h}=\tanh(\mathbf{x}W_{\mathrm{in}})\) and
\(\mathbf{a}=(\pi/2)\tanh(\mathbf{h}W_{\mathrm{out}})\).
A constant input feature supplies a hidden-layer bias. Input and output sizes
remain fixed across all network sizes. Genotypes explicitly encode dimensions
and the two weight matrices. For $I$ input features (including the constant
feature), $H$ hidden neurons and $O$ actuators, the connection count is
$P=IH+HO$. There is no independent output-bias vector in this proposed encoding.

\subsection{Conditions}
\begin{table}[t]
\caption{Proposed conditions. Replace parameters if the pilot changes them.}
\label{tab:conditions}
\centering
\begin{tabular}{lll}
\toprule
Condition & Hidden neurons & Search\\
\midrule
Fixed small & 2, fixed & Weights\\
Fixed large & 8, fixed & Weights\\
Growing & 2 initially; cap 8 & Weights + node addition\\
Random small & 2, fixed & Independent sampling\\
\bottomrule
\end{tabular}
\end{table}

The small fixed network controls for the growing network's starting capacity.
The large fixed network tests whether additional capacity available from the
start accounts for any improvement. Random search is a baseline using the small
network representation and the same initialization distribution and evaluation
budget. It does not separately characterize random search of the larger network.

\subsection{Shared evolutionary algorithm}
\todo{Person 1: replace the proposed description with the implemented pipeline.
Specify mu/lambda, selection, tie handling, initialization distribution,
mutation, evaluation caching, minimization and ARIEL EC integration.}

The proposed EA uses tournament parent selection and $(\mu+\lambda)$ elitist
survivor selection. Weight mutation independently perturbs each weight with
probability $1/P$, adding Gaussian noise with standard deviation $\sigma$.
The expected number of perturbed weights is therefore one, regardless of network
size. Identical weight mutation and selection are used in all EA conditions.
No crossover or speciation is included. Both NumPy and ARIEL random generators
are seeded for each independent run.

\subsection{Topology-growth operator}
\todo{Person 2: confirm the operation order, cap behavior and validation. Describe
exactly what changes and what remains fixed. Do not call this full NEAT.}

After weight mutation, the proposed growing condition adds one hidden neuron
with probability $p_{\mathrm{add}}$, unless the cap has been reached. The new
incoming weights are sampled from the initialization distribution; its outgoing
weights start at zero. Existing weights are preserved. Node addition alone
therefore preserves the pre-addition controller function, while later weight
mutations can give the new neuron an effect. Zero-weight connections count as
connections. Ties are randomized rather than resolved in favor of smaller
networks, allowing neutral additions to survive.

This operator expands the width of one dense hidden layer. It does not add
layers, remove nodes, introduce recurrent connections, or implement historical
markings. Performance conclusions apply to this restricted form of growth.

\subsection{Budget, repeats and analysis}
\todo{Person 3, checked by both coders: insert the final parameter values,
independent seeds, pilot choices, hardware and evaluation-budget accounting.}

\begin{table}[t]
\caption{Provisional settings to freeze after the pilot. These are not measured results.}
\label{tab:settings}
\centering
\begin{tabular}{lr}
\toprule
Parameter & Proposed value\\
\midrule
Population $\mu$ / offspring $\lambda$ & 20 / 20\\
Tournament size & 3\\
Initialization SD & 0.5\\
Weight-mutation SD $\sigma$ & 0.1\\
Expected mutated weights & 1\\
Node-addition probability & 0.1\\
Evaluation budget per run & 1,020\\
Independent final seeds & 0, 1, 2, 3, 4\\
Simulation duration & 15 s\\
\bottomrule
\end{tabular}
\end{table}

Initialization evaluations are included in the budget. Evaluated survivors are
cached; each newly simulated offspring consumes one evaluation. Random search
uses the same total number of simulations and is logged in batches matching the
EA's evaluation increments. Pilot seeds are kept separate from final seeds.
Pilot runtime and plateau behavior determine a common final budget; all final
runs use that cap instead of separate early stopping rules.

For each condition, the analysis reports the mean and sample standard deviation
of final best-so-far target distance across independent runs. Fitness curves
summarize each run's best-so-far distance and then aggregate across seeds.
Within-population spread is not used as a substitute for across-run variation.
Both generation/batch and evaluation-count plots are retained. Population mean
hidden-neuron and connection counts describe how topology develops.
\todo{Explain any deviation from the proposal. With five repeats, prioritize
effect sizes and variation; do not claim significance without a justified test.}

\section{Results and Discussion}
\subsection{Target-reaching performance}
\todo{Person 3: insert actual per-condition means and sample SDs, state the
budget, and compare growing against BOTH fixed conditions and the baseline.
No results currently exist in this template.}

\begin{figure}[t]
\centering
\figureplaceholder{fitness_generations.pdf}{Fitness across generations / random-search batches}
\caption{Replace with observed mean best-so-far distance and $\pm$ one sample
standard deviation across independent runs. State $n$, budget and whether any
runs failed. Lower distance is better.}
\label{fig:fitness}
\end{figure}

\begin{table}[t]
\caption{Final best-so-far distance. Dashes are empty placeholders, not zeros.}
\label{tab:results}
\centering
\begin{tabular}{lrrr}
\toprule
Condition & Runs & Mean (m) & SD (m)\\
\midrule
Fixed small & --- & --- & ---\\
Fixed large & --- & --- & ---\\
Growing & --- & --- & ---\\
Random small & --- & --- & ---\\
\bottomrule
\end{tabular}
\end{table}

\todo{Use the evaluation-count plot as a second figure if it adds information;
otherwise the generation plot already has identical evaluation increments.
Show individual final seed values when useful and affordable within six pages.}

\subsection{Topology development}
\begin{figure}[t]
\centering
\figureplaceholder{topology_growth.pdf}{Hidden-neuron and connection counts across generations}
\caption{Replace with measured population mean sizes, aggregated across runs.
State whether growth reaches the cap and whether added capacity improves fitness.}
\label{fig:topology}
\end{figure}

\todo{Person 2 supplies observations: do neurons survive after addition? Do new
outgoing weights become useful? Does the best controller grow, or just the
population? Distinguish correlation from evidence of a causal growth benefit.}

\subsection{Interpretation and limitations}
\todo{Answer the question using results, including negative or mixed outcomes.
If growing beats small but matches large, additional capacity is a plausible
explanation; avoid claiming staged growth is uniquely responsible. Discuss
limited seed count, finite budget, one body/world, one-layer topology, hard cap,
function-preserving initialization, and omitted NEAT mechanisms. Report whether
curves plateau or the study remains budget-limited. Equal simulation counts do
not guarantee equal computation time.}

\section{Conclusions}
\todo{Person 3: give a direct, qualified answer to the research question, supported
by the final-distance and growth evidence. State the scope of the conclusion and
one focused next experiment. Do not introduce new unreported results.}

% Start references on a separate page so content pages can be counted clearly.
\clearpage
\begin{thebibliography}{9}
\bibitem{stanley2002}
Kenneth O. Stanley and Risto Miikkulainen. 2002.
Evolving Neural Networks through Augmenting Topologies.
\emph{Evolutionary Computation} 10, 2, 99--127.
\url{https://nn.cs.utexas.edu/downloads/papers/stanley.ec02.pdf}

\bibitem{lecture2026}
Karine Miras. 2026. \emph{Neuroevolution}. Evolutionary Computing lecture slides,
\texttt{EC\_NeuroEvolution\_2026\_v2.pdf}.

\bibitem{assignment2026}
Evolutionary Computing course. 2026.
\emph{Standard Assignment -- Task II}. Assignment handout and
\texttt{A2\_template\_2026.py}. Course materials.

\bibitem{ariel}
ARIEL contributors. \emph{ARIEL documentation}.
\url{https://ci-group.github.io/ariel/}.
% Add the actual code version/commit in Methods once confirmed by both coders.
\end{thebibliography}
\end{document}
